\chapter{Electronic Structure}
\label{ch:electronic}

\menu{Simulation}\menusep\menu{Electronic Structure\dots} computes a band
structure along a high-symmetry $k$-path and, where the backend supports
it, an element- and orbital-resolved projected density of states.

\begin{calwarn}
The GPAW route needs a \textbf{baseline ground state}: a completed
\menu{Simulation}\menusep\menu{Single-point Calculation\dots} run with the GPAW
calculator, which writes \file{single\_point.gpw} into its job directory. The
bands run loads that converged density and evaluates non-self-consistently at
fixed density --- it never re-runs the SCF. Without a baseline the wizard
refuses to open and tells you so. Chapter~\ref{ch:tutorial-silicon} walks
through the whole sequence on silicon.
\end{calwarn}

\section{Setting up the calculation}

\begin{description}
  \item[\ui{Baseline SCF density}] The completed single point whose
    \file{single\_point.gpw} supplies the charge density. Inheriting it rather
    than re-converging is what makes the band structure attributable to a
    specific, inspected ground state.
  \item[\ui{Backend}] \ui{Free electrons (ASE --- bands only)},
    \ui{GPAW (DFT, bands + PDOS)} or
    \ui{Quantum ESPRESSO (needs pw.x + pseudos)}. The GPAW entry is
    disabled when the package is not importable in the selected
    environment.
  \item[\ui{k-path}] Pre-filled with ASE's suggestion for the structure's
    Bravais lattice, e.g.\ \texttt{GXWKGLUWLK,UX}. Edit it freely; leave it
    empty to accept the suggestion.
  \item[\ui{k-points}] Total points along the whole path, default 80.
  \item[\ui{Valence electrons}] Electrons per cell, free-electron backend
    only.
  \item[\ui{PW cutoff}] Plane-wave cutoff, default 340\,eV.
  \item[\ui{SCF k-grid}] Monkhorst-Pack $n^3$ grid for the self-consistent
    step, default 4.
  \item[\ui{Compute element/orbital PDOS (GPAW backend)}] Adds the
    projected DOS.
  \item[\ui{Environment}] The Conda environment or interpreter that runs
    the job (Section~\ref{sec:exec-env}).
\end{description}

\begin{caltip}
The free-electron backend needs nothing beyond ASE and finishes in seconds.
It computes empty-lattice bands --- the exact free-electron dispersion
folded into your Brillouin zone --- which is genuinely useful for checking
that a $k$-path is the one you meant, and for teaching band folding, before
committing a DFT run.
\end{caltip}

\begin{calwarn}
Real bands need a real solver: \texttt{gpaw} installed in the selected
environment, or \texttt{pw.x} plus pseudopotentials for the Quantum
ESPRESSO route. The QE script is a scaffold with an \texttt{EDIT ME} line
for your pseudopotential set.
\end{calwarn}

The job runs like any other --- console, progress markers, process-manager
entry --- and writes \file{bands.json} and, when computed,
\file{pdos.json} into the job directory. The viewer opens automatically on
completion, and can be reopened any time with \ui{Load Result} in the
process manager.

\section{Reading the plots}

The viewer shows the band structure and, when present, the PDOS side by
side on a shared energy axis.

\subsection{Band structure}

Energy versus $k$-path distance, with vertical lines and labels at every
high-symmetry point (ASE's \texttt{G} is rendered as $\Gamma$). Spin
channels are drawn in different colours. A dashed horizontal line marks the
reference energy at $E - E_\text{ref} = 0$.

\subsection{Projected density of states}

Density of states versus the same energy axis, one curve per element and
orbital projection --- \texttt{Si s}, \texttt{Si p}, \texttt{O p} and so on
--- accumulated over all atoms of each species and colour-keyed to the
legend.

\section{Controls and export}

\begin{description}
  \item[\ui{Fermi level}] The reference energy, pre-filled from the
    calculation. Plots show $E - E_\text{ref}$, so shifting it moves the
    zero line --- useful for aligning against a reference calculation or a
    measured spectrum.
  \item[\ui{E min}, \ui{E max}] The energy window, defaults $-10$ and
    $+10$\,eV.
  \item[\ui{Projections}] A checklist of the PDOS curves. Uncheck the ones
    that clutter the picture; the vertical scale re-normalises to what
    remains visible.
  \item[\ui{Export Bands\ldots}] CSV or \file{.dat}: one row per $k$-point,
    with the path distance followed by every band (and spin channel).
  \item[\ui{Export PDOS\ldots}] CSV or \file{.dat}: one row per energy, one
    column per projection.
\end{description}

\begin{calnote}
Exported bands use the same $k$-path distance as the $x$ axis, so the
columns drop straight into an external plotting script and reproduce the
figure you see.
\end{calnote}

\section{X-ray absorption spectroscopy (XAS)}
\label{sec:xas}

\menu{Electronics}\menusep\menu{X-ray Absorption Spectroscopy (XAS)\dots}
computes a core-level absorption spectrum with GPAW, following the GPAW XAS
tutorial.

An XAS transition starts in a \emph{core} level of one atom, so it needs a PAW
dataset with a hole in that level --- and none ships with GPAW. The generated
script therefore runs three stages rather than one: it builds the core-hole
dataset for the absorbing element, converges a ground state that uses it on the
absorbing \emph{atom}, and evaluates the spectrum from those wavefunctions. The
dataset is written into the job directory and found through
\texttt{setup\_paths}, so nothing is installed into your GPAW distribution and
the run reproduces elsewhere.

\subsection{Absorbing site}

Pick the \ui{Element} and then the \ui{Absorbing atom} --- one atom, not all
atoms of that species. Giving the core-hole setup to every one of them would
model a solid in which all are excited simultaneously, which is not what an
absorption measurement does. Inequivalent sites each need their own run, and
the measured spectrum is their sum; the wizard says so when the structure has
more than one candidate.

\ui{Edge} selects the level: K (1s) is what almost every XAS measurement means,
while the L edges (2s, 2p) matter for transition metals.

\subsection{Core hole}

\begin{description}
  \item[\ui{Half hole}] The transition-potential approximation --- 0.5
    electrons removed. One calculation gives the whole spectrum, because the
    final-state relaxation is averaged between initial and final states. This
    is the tutorial's choice and what most published XAS is.
  \item[\ui{Full hole}] The excited final state proper. Better for the first
    resonance, worse for the rest of the spectrum, and it needs the
    delta-Kohn-Sham shift to sit on an absolute energy scale.
  \item[\ui{No hole}] The unperturbed ground state --- what the spectrum would
    be if the core hole did not pull the excited states down. Worth computing
    to see how large that effect is.
\end{description}

\ui{Unoccupied bands} sets how many states above the occupied ones are
converged. The spectrum is a sum over them, so too few truncates it --- visibly,
as a spectrum that stops rather than decays.

\subsection{Broadening and the energy scale}

\ui{Broadening (FWHM)} is a uniform Gaussian on every transition.
\ui{Broaden more at higher energy} ramps it linearly above the edge, which is
the physical behaviour (core-hole lifetime plus final-state broadening) and what
makes the result comparable with a measurement whose peaks wash out with energy.
The default ramp is the tutorial's, chosen for the oxygen K edge --- move it to
your own.

GPAW produces the spectrum on a \emph{relative} scale whose zero is the first
unoccupied state, which is not where a measurement puts it.
\ui{Compute the absolute edge position} adds two total-energy calculations
(delta-Kohn-Sham) that give the shift, roughly doubling the cost; alternatively
type a known \ui{Edge energy} directly.

\begin{calnote}
\texttt{gpaw.xas} runs on GPAW's \emph{legacy} engine only. This is the one
script Calango generates that clears \texttt{GPAW\_NEW} and asks for
\texttt{legacy\_gpaw=True} --- without it the run stops with ``New-GPAW not
supported''.
\end{calnote}

\subsection{Results}

The viewer opens automatically when the run finishes and plots the isotropic
spectrum --- the average over the three Cartesian polarizations, which is what a
powder or solution measurement sees. \ui{Show the x, y and z polarizations}
adds them individually; for an oriented sample the difference between them is
the result. \ui{Show individual transitions} overlays the discrete lines the
broadened curve is built from, which is how you tell a genuine shoulder from two
peaks the broadening merged.

The header states which energy scale the axis is on, because a relative spectrum
plotted as though it were absolute is wrong by hundreds of eV and nothing about
the curve reveals it.

\section{Unfolding a relaxed supercell}
\label{sec:unfoldrelaxed}

Unfolding needs the supercell to be an exact integer multiple of the primitive
cell: the projection is defined against that mapping. A \emph{relaxed} doped
supercell is not --- substituting an atom and letting the cell relax moves its
lattice constant by a per cent or so, and the pristine host is no longer
$1/n$ of it.

The \ui{Effective Bands} wizard handles that two ways.

\ui{Tolerance} is how far $M$ may sit from integers before the two cells are
called incommensurate. The default (0.02) has room for a relaxation; an
unrelated primitive cell misses by 0.1 or more, so the check still catches the
mistake it exists for.

\ui{Force commensurability by rescaling the unit cell} rebuilds the primitive
cell as $P = M^{-1} S$, making the relation exact. For a relaxed supercell this
is not a fudge: the supercell \emph{is} $M \times$ something, just not
$M \times$ the pristine host any more, and this recovers the host cell it
actually relaxed into. Only the lattice is rebuilt --- the basis rides along at
the same fractional coordinates, so every atom stays on its site.

The wizard reports the resulting strain. A couple of per cent is the
relaxation; much more means the wrong primitive cell was selected, and it says
so.

\begin{calnote}
Worked example --- Au$_7$Pt. A primitive Au fcc cell ($a = 4.078$\,\AA) and a
$2\times2\times2$ supercell with one Au replaced by Pt, relaxed with GPAW,
comes back at $a = 4.1088$\,\AA: a 0.755\% expansion. The deduced
$M = 2\mathbf{I}$ is right, but its residual is 0.015 --- rejected outright by
a pristine-cell tolerance, accepted by the default one. Forcing rebuilds the
primitive at 4.1088\,\AA{} and the residual falls to $5\times10^{-17}$.
\end{calnote}
